\SourceLeaf{494}%
sind die Theodolite, wie sie der englische Instrumentenmacher \textsc{Troughton} anfertigt, am bequemsten und brauchbarsten. Prof. \textsc{Schumacher} war so freundlich, mir seinen für diesen Zweck zum Einkaufspreis abzutreten. Er hofft, bis zu dem Zeitpunkt, zu dem er selbst in diesem Jahr wieder einen braucht, einen anderen aus England zu erhalten. — Manche Winkelmessungen werden sich kaum anders als nachts mit Argandlampen und parabolischen Reflektoren anstellen lassen, die über sehr große Entfernungen sichtbar gemacht werden können. Damit es bei Bedarf daran nicht fehlt, habe ich mehrere solche Reflektoren bei einem sehr geschickten Instrumentenmacher bestellt, dem Hofmechaniker \textsc{Körner} in Jena, der auch ähnliche Arbeiten für die preußischen Messungen geliefert hat. Er hat versprochen, sie bis Ostern fertigzustellen. Mehrere andere im Voraus besprochene Apparate lasse ich hier unerwähnt.

Nach all diesen Vorkehrungen glaube ich nun, die Arbeiten im kommenden Frühjahr beginnen zu können. Wenn die Instrumentenmacher alle ihre Versprechen gehörig einhalten, kann das Werk dann weiterhin rasch voranschreiten. In jedem Fall werden aber die vielfältigen Vorarbeiten, die Bereisung der Gegenden, das Aussuchen der Stationen, der Bau von Signaltürmen, die vorläufigen Messungen und dergleichen — (den zuvor erwähnten Theodoliten von \textsc{Troughton} habe ich bereits) — zunächst mehrere Monate ausfüllen. Und selbst wenn der ungünstige Fall einträte, dass der erwartete größere Theodolit von \textsc{Reichenbach} auch dann noch nicht geliefert wäre, wenn die Vorarbeiten so weit fortgeschritten sind, dass die Messungen selbst beginnen können, müsste die Arbeit nicht zum Stillstand kommen. Nötigenfalls könnte ich mit den Winkelmessungen anfangen, für die die übrigen Hilfsmittel der Sternwarte ausreichen. Ja selbst im schlimmsten Fall, der hoffentlich nicht eintritt, dass die Lieferung des Theodoliten sich noch länger verzögert, würde ich nur die Reihenfolge der Arbeiten ändern. Ich würde den Rest des Sommers für die astronomischen Arbeiten mit dem Zenitsektor hier in Göttingen verwenden, die sonst, wenn der ganze Sommer der Triangulation gewidmet werden kann, einer späteren Zeit vorbehalten bleiben.

Göttingen, den 20. März 1821.
\begin{flushright}C. F. \textsc{Gauss}.\end{flushright}

\par\bigskip\Needspace{13\baselineskip}
\pdfbookmark[0]{Druckfehler.}{druckfehler}
\begin{center}\large DRUCKFEHLER.\end{center}
\par\noindent Auf Seite 296 statt
\begin{center}
\begin{tabular}{c|c|c|c}
\(Q+q\)&\(P+p\)&\(\log m\)&\(k\)\\
&&\(+\)&\\
\(52^{0}\ 30'\)&\(52^{0}\ 32'\ 1''78428\)&&\(0''006\)
\end{tabular}
\end{center}
\noindent lies
\begin{center}
\begin{tabular}{c|c|c|c}
\(Q+q\)&\(P+p\)&\(\log m\)&\(k\)\\
&&\(-\)&\\
\(52^{0}\ 30'\)&\(52^{0}\ 32'\ 1''78428\)&&\(0''006\)
\end{tabular}
\end{center}
